\section{图论}

\subsection{Dijkstra 算法}
\begin{lstlisting}[caption={}]
std::vector<i64> dist(n + 1),vis(n + 1);
auto dijkstra = [&](int st) -> void{
	std::fill(dist.begin() + 1,dist.end(),1e10);
	std::fill(vis.begin() + 1,vis.end(),0);
	std::priority_queue<std::array<i64,2>,std::vector<std::array<i64,2>>,std::greater<std::array<i64,2>>> pq;
	dist[st] = 0;
	pq.push({0,st});
	while(!pq.empty()){
		auto [d,u] = pq.top();
		pq.pop();
		if(vis[u]) continue;
		vis[u] = 1;
		for(auto [v,w] : edge[u]){
			if(dist[v] > d + w){
				dist[v] = d + w;
				pq.push({d + w,v});
			}
		}
	}
};
dijkstra(s);
\end{lstlisting}


\subsection{Floyd算法}
\begin{lstlisting}[caption={}]
int n,m,k,x,y,z;
int dist[N][M];
//dist[i][j] 表示点i到点j的最短距离 ,k为中间点 
void floyd(){
	for(k = 1;k <= n;k++){   //枚举中间点 
		for(int i = 1;i <= n;i++){ //起始点 
			for(int j = 1;j <= n;j++){ //目的地 
				dist[i][j] = min(dist[i][j] , dist[i][k] + dist[k][j]);
			}
		}
	}
} 
\end{lstlisting}

\subsection{二分图最大匹配(匈牙利算法)}
\begin{lstlisting}[caption={}]
#include<bits/stdc++.h>
using namespace std;

typedef long long ll;
const int N = 5e2 + 10;

ll n,m,fri[N],vis[N];
vector<ll> e[N];

bool match(ll x,ll t){
    if(vis[x] == t) return false; //已经访问过 
	vis[x] = t; //记录状态为访问过
	for(auto i : e[x]){
		if(fri[i] == 0 || match(fri[i],t)){ //如果暂无匹配，或原来匹配的左侧元素可以找到新的匹配
			fri[i] = x; //当前左侧元素成为当前右侧元素的新匹配
			return true; //返回匹配成功
		}
	}
	return false; //循环结束，仍未找到匹配，返回匹配失败
}

int main(){
	ios::sync_with_stdio(0),cin.tie(0),cout.tie(0);
	ll k,a,b,cnt = 0;
	cin >> n >> m >> k;
	while(k--){
		cin >> a >> b;
		e[a].push_back(b);
	}
	for(int i = 1;i <= n;i++)
		if(match(i,i)) cnt++;
	cout << cnt << "\n";
	return 0;
}
\end{lstlisting}


\subsection{差分约束}
给出一组包含 $m$ 个不等式,有 $n$ 个未知数的形如:
$$ \begin{cases} u_1-v_1\leq w_1 \\u_2-v_2 \leq w_2 \\ \cdots\\ u_m -v_m\leq w_m\end{cases}$$
的不等式组，求任意一组满足这个不等式组的解。$\sf SPFA$ 解,$\mathcal O(nm)$ 。
\begin{lstlisting}[caption={}]
#include<bits/stdc++.h>
#define all(x) (x).begin(),(x).end()
using namespace std;

typedef int ll;
typedef pair<ll,ll> PII;
const int N = 5e3 + 10,M = 5e3 + 10,mod = 1e9 + 7;

int head[N];
struct edge{
    int next,to, w;
};
vector<edge> edges;

void add_edge(int u, int v, int w){
    edges.push_back({head[u],v,w});
    head[u] = edges.size() - 1;
}
bool inqueue[N];
int cnt[N], dis[N];
queue<int> Q;
bool SPFA(int s, int n){
    memset(dis, 127, sizeof(dis));
    dis[s] = 0;
    Q.push(s);
    while (!Q.empty()){
        int p = Q.front();
        if (cnt[p] > n)
            return false;
        Q.pop();
        inqueue[p] = false;
        for (int eg = head[p]; eg != -1; eg = edges[eg].next){
            int to = edges[eg].to;
            if (edges[eg].w + dis[p] < dis[to]){
//          if (edges[eg].w + dis[p] > dis[to])
                dis[to] = edges[eg].w + dis[p];
                if (!inqueue[to]){
                    Q.push(to);
                    inqueue[to] = true;
                    cnt[to]++;
                }
            }
        }
    }
    return true;
}

int main(){
    ios::sync_with_stdio(0),cin.tie(0),cout.tie(0);
    ll n,m,x,y,w;
    cin >> n >> m;
    memset(head,-1,sizeof(head));
    for (int i = 0; i < m; ++i){
        cin >> x >> y >> w;
        add_edge(y, x, w);
    }
    for (int i = 1; i <= n; ++i) add_edge(0, i, 0);
    if (SPFA(0, n)){ //SPFA判断负环，无解
        for (int i = 1; i <= n; ++i)
           cout << dis[i] << " ";
    }
    else cout << "NO\n";
	return 0;
}
\end{lstlisting}

\subsection{2-SAT}
基于 tarjan 缩点，时间复杂度为 $\mathcal O(N+M)$。
\begin{lstlisting}[caption={}]
#include<bits/stdc++.h>
using namespace std;
typedef long long ll;

struct SCC{
    // st：是否在栈中  scc：每个点所属的强连通分量编号  cscc：强连通分量的数量
    vector<ll> dfn,low,stk,scc;
    vector<bool> st;
    vector<vector<ll>> e;
    ll cscc,N,cnt = 0;
    SCC(int n){
        dfn.resize(n + 1,0);
        low.resize(n + 1,0);
        st.resize(n + 1,false);
        scc.resize(n + 1);
        e.resize(n + 1);
        cscc = 0;
        N = n;
    }

    void add(ll u,ll v){
        e[u].push_back(v);
    }

    void tarjan(ll p){
        low[p] = dfn[p] = ++cnt;
        st[p] = true;
        stk.push_back(p);
        for(auto& q : e[p]){
            if(!dfn[q]){
                tarjan(q);
                low[p] = min(low[q],low[p]);
            }
            else if(st[q]){
                low[p] = min(low[p],dfn[q]);
            }
        }
        if(low[p] == dfn[p]){
            int top;
            cscc++;
            do{
                top = stk.back();
                stk.pop_back();
                st[top] = false;
                scc[top] = cscc;
            }while(top != p);
        }
    }

    void work(){
        for(int i = 1;i <= N;i++)
            if(!dfn[i])
                tarjan(i);
    }
};
//scc编号的顺序符合拓扑序（编号越大的点拓扑序越靠前）

int main(){
	ios::sync_with_stdio(0),cin.tie(0),cout.tie(0);
    ll n,m;
    cin >> n >> m;
    SCC sc(2 * n + 1);
    for(int i = 1;i <= m;i++){
        ll a,ta,b,tb;
        cin >> a >> ta >> b >> tb;
        sc.add((2 * a + ta) ^ 1,2 * b + tb);
        sc.add((2 * b + tb) ^ 1,2 * a + ta);
    }
    sc.work();
    bool flag = true;
    vector<ll> ans(n + 1);
    for(int i = 1;i <= n;i++){
        if(sc.scc[2 * i] == sc.scc[2 * i + 1]){
            flag = false;
            break;
        }
        else if(sc.scc[i * 2] > sc.scc[i * 2 + 1])
            ans[i] = 1;
        else ans[i] = 0;
    }
    if(flag){
        cout << "POSSIBLE\n";
        for(int i = 1;i <= n;i++){
            cout << ans[i] << " ";
        }
    }
    else cout << "IMPOSSIBLE\n";
	return 0;
}
\end{lstlisting}

\subsection{割边}
\begin{lstlisting}[caption={}]
//bri为割边
struct CutBridge{
    vector<ll> dfn,low,fa;
    vector<vector<ll>> edge;
    int n,m,cnt = 0;
    vector<PII> bri;  

    CutBridge(int n,int m){
        edge.resize(n + 2),dfn.resize(n + 1,0),low.resize(n + 1,0);
        fa.resize(n + 1,0);
        this->n = n,this->m = m;
    }

    void add(int a,int b){
        edge[b].push_back(a);
        edge[a].push_back(b);
    }

    void tarjan(int p){
        low[p] = dfn[p] = ++cnt;
        for (auto to : edge[p]){
            if (!dfn[to]){
                fa[to] = p; // 记录父节点
                tarjan(to);
                low[p] = min(low[p], low[to]);
                if (low[to] > dfn[p])
                    bri.emplace_back(p, to);
            }
            else if (fa[p] != to) // 排除父节点
                low[p] = min(low[p], dfn[to]);
        }
    }

    void work(){
        for(int i = 1;i <= n;i++)
            if(!dfn[i])
                tarjan(i);
    }
};
\end{lstlisting}

\subsection{割点}
\begin{lstlisting}[caption={}]
vector<ll> edge[N],cut;//cut 存储所有割点
ll dfn[N],low[N],cnt = 0;
void tarjan(ll p,bool root = true){
    ll tot = 0;
    cnt++;
    dfn[p] = low[p] = cnt;
    for(auto& i : edge[p]){
        if(!dfn[i]){
            tarjan(i,false);
            low[p] = min(low[p],low[i]);
            tot += (low[i] >= dfn[p]);
        }
        else low[p] = min(low[p],dfn[i]);
    }
    if(tot > root)
        cut.push_back(p);
}
\end{lstlisting}

\subsection{缩点}
\begin{lstlisting}[caption={}]
struct SCC{
    // st：是否在栈中  scc：每个点所属的强连通分量编号  cscc：强连通分量的数量
    vector<ll> dfn,low,stk,scc;
    vector<bool> st;
    vector<vector<ll>> e;
    ll cscc,N,cnt = 0;
    SCC(int n){
        dfn.resize(n + 1,0);
        low.resize(n + 1,0);
        st.resize(n + 1,false);
        scc.resize(n + 1);
        e.resize(n + 1);
        cscc = 0;
        N = n;
    }

    void add(ll u,ll v){
        e[u].push_back(v);
    }

    void tarjan(ll p){
        low[p] = dfn[p] = ++cnt;
        st[p] = true;
        stk.push_back(p);
        for(auto& q : e[p]){
            if(!dfn[q]){
                tarjan(q);
                low[p] = min(low[q],low[p]);
            }
            else if(st[q]){
                low[p] = min(low[p],dfn[q]);
            }
        }
        if(low[p] == dfn[p]){
            int top;
            cscc++;
            do{
                top = stk.back();
                stk.pop_back();
                st[top] = false;
                scc[top] = cscc;
            }while(top != p);
        }
    }

    void work(){
        for(int i = 1;i <= N;i++)
            if(!dfn[i])
                tarjan(i);
    }
};
//scc编号的顺序符合拓扑序（编号越大的点拓扑序越靠前）

int main(){
	ios::sync_with_stdio(0),cin.tie(0),cout.tie(0);
    ll n,m,ans = 0;
    cin >> n >> m;
    SCC a(n);
    vector<ll> w(n + 1),nw;
    for(int i = 1;i <= n;i++) cin >> w[i];
    for(int i = 1;i <= m;i++){
        ll u,v;
        cin >> u >> v;
        a.add(u,v);
    }
    a.work();
    nw.resize(a.cscc + 1,0);
    vector<vector<ll>> edge(a.cscc + 1);
    vector<ll> dp(a.cscc + 1,0);
    for(int i = 1;i <= n;i++){
        nw[a.scc[i]] += w[i];
        dp[a.scc[i]] += w[i];
        ans = max(ans,dp[a.scc[i]]);
        for(auto& j : a.e[i]){
            if(a.scc[i] != a.scc[j]){
                edge[a.scc[i]].push_back(a.scc[j]);
            }
        }
    }
    for(int i = a.cscc;i >= 1;i--){
        for(auto& j : edge[i]){
            dp[j] = max(dp[j],dp[i] + nw[j]);
            ans = max(ans,dp[j]);
        }
    }
    cout << ans << "\n";
	return 0;
}
\end{lstlisting}




\subsection{边双连通分量}
\begin{lstlisting}[caption={}]
struct EBCC{
    // st：是否在栈中  ebcc：每个点所属的边双连通分量编号  cebcc：边双连通分量的数量
    std::vector<i64> dfn,low,stk,ebcc;
    std::vector<bool> st;
    std::vector<std::vector<std::array<i64,2>>> e;
    i64 cebcc,N,cnt,edge_cnt;
    EBCC(int n): dfn(n + 1),low(n + 1),st(n + 1),
                 ebcc(n + 1),e(n + 1),cebcc(0),N(n),
				 cnt(0),edge_cnt(1){}
    void add(i64 u,i64 v){
        e[u].push_back({v,edge_cnt << 1});
        e[v].push_back({u,edge_cnt << 1 | 1});
		edge_cnt++;
    }
    void tarjan(int p, int las){
		low[p] = dfn[p] = ++cnt;
		st[p] = true;
        stk.push_back(p);
		for(auto [q,idx]: e[p]){
			if(idx == (las ^ 1)) continue;
			if(!dfn[q]){
				tarjan(q, idx);
				low[p] = std::min(low[p], low[q]);
			}else low[p] = std::min(low[p], dfn[q]);
		}
		if(low[p] == dfn[p]){
            int top;
            cebcc++;
            do{
                top = stk.back();
                stk.pop_back();
                st[top] = false;
                ebcc[top] = cebcc;
            }while(top != p);
        }
	}
    void work(){
        for(int i = 1;i <= N;i++)
            if(!dfn[i])
                tarjan(i,0);
    }
    // 获取缩点后的树（边双连通分量缩点后的树）
    std::vector<std::vector<i64>> get_tree(){
        std::vector<std::vector<i64>> tree(cebcc + 1);
        for(int u = 1; u <= N; u++){
            for(auto& [v,_] : e[u]){
                if(ebcc[u] != ebcc[v]){
                    tree[ebcc[u]].push_back(ebcc[v]);
                }
            }
        }
        for(int i = 1; i <= cebcc; i++){
            std::sort(tree[i].begin(), tree[i].end());
            tree[i].erase(std::unique(tree[i].begin(), tree[i].end()), tree[i].end());
        }
        return tree;
    }
};
\end{lstlisting}


\subsection{欧拉路径}
\subsubsection{有向图的欧拉路径}
有向图欧拉路径存在:恰有一个点出度比入度多1(为起点);恰有一个点入度比出度多 
$1$(为终点);恰有$N - 2$ 个点入度均等于出度。如果是欧拉回路，则上方起点与终点的条件不存在，
全部点均要满足最后一个条件。
\begin{lstlisting}[caption={}]
//有向图欧拉路径求解（字典序最小）
vector<int> ans;
auto dfs = [&](auto self, int x) -> void {
    while (ver[x].size()) {
        int net = *ver[x].begin();
        ver[x].erase(ver[x].begin());
        self(self, net);
    }
    ans.push_back(x);
};
dfs(dfs, s);
reverse(ans.begin(), ans.end());
for (auto it : ans) {
    cout << it << " ";
}
\end{lstlisting}

\subsubsection{无向图的欧拉路径}
无向图欧拉路径存在：恰有两个点度数为奇数(为起点与终点);恰有 $N - 2$个点度数为偶数。
\begin{lstlisting}[caption={}]
//无向图欧拉路径求解
auto dfs = [&](auto self, int x) -> void {
    while (ver[x].size()) {
        int net = *ver[x].begin();
        ver[x].erase(ver[x].find(net));
        ver[net].erase(ver[net].find(x));
        cout << x << " " << net << endl;
        self(self, net);
    }
};
dfs(dfs, s);
\end{lstlisting}